\section{Related Work}
\label{sec:realted_work}

\subsection{Hallucination Detection Benchmarks}
Hallucination detection aims to develop a framework or a model to automatically distinguish between hallucinated and factual content~\cite{ li2025hd,ravichander2025halogen,qin2025learning,zhang2025icr,zhang2025prompt}. As shown in Table~\ref{dataset_comparison}, a line of work assesses the model’s factuality reasoning over diverse domains such as world knowledge~\cite{li2023halueval, wei2024long,bang2025hallulens}, science~\cite{zhao2023felm}, and math~\cite{zhao2023felm}. Moreover, RAGTruth~\cite{niu2024ragtruth} demonstrates that popular LLMs continue to hallucinate across tasks even with retrieval-augmented generation. To diagnose tool-use hallucinations, ToolBH~\cite{zhang2024toolbehonest} collects 700 tool-call samples to perform solvability detection, solution planning, and missing-tool analysis. Different from prior works confined to binary judgment in single-turn responses, we introduce the first benchmark for automated hallucination attribution within multi-step agent trajectories.

\subsection{LLM-based Agents}
LLM-based agents~\cite{yao2022react,wang2024executable} have showcased extraordinary capabilities in automating tasks across various fields. This growing capability is largely driven by emergent behaviors that arise during chain-of-thought~\cite{wei2022chain}, in-context learning~\cite{brown2020language}, and instruction following~\cite{longpre2023flan}. To extend agents' capabilities beyond their internal knowledge, function calling~\cite{schick2023toolformer,patil2024gorilla,patilberkeley} has been proposed, enabling agents to interact with external tools and APIs in multi-step workflows.

Moreover, individual agents, each serving a specialized role, can be composed into multi-agent systems to solve complex tasks~\cite{hong2023metagpt, wang2024openhands}. Early multi-agents elicit stepwise reasoning through structured debate~\cite{liang2024encouraging} or role-play dialogue~\cite{li2023camel}. Recent works~\cite{fourney2024magentic, hu2025owl} introduce central orchestrators that assign tasks to specialized agents. Beyond inter-agent interaction, agents are motivated to proactively seek human feedback to improve their decision-making~\cite{feng2024large}. Despite this progress, hallucinations persist across operational stages in agent workflows~\cite{lin2025llm}, including planning, retrieval, reasoning, tool use, and human interaction, underscoring the need for reliability assessment.






